% Source: https://artofproblemsolving.com/community/c3h1404200_math_kangaroo_2015_sample_questions__levels_7__8_prob_14
% Stable identity: aops-topic-1404200-q14
"$15$ litres per square metre" means that every square metre of ground (or water surface) received $15$ litres. Since $1$ litre $=1000\text{ cm}^3$ and $1\text{ m}^2=100\text{ cm}\times100\text{ cm}=10000\text{ cm}^2$, a column of rain standing on $1\text{ m}^2$ has height
\[
\dfrac{15000\text{ cm}^3}{10000\text{ cm}^2}=1.5\text{ cm}.
\]
A pool of surface area $A$ square metres receives $15A$ litres spread over the same area $A$, so its level rises by the same $1.5$ cm whatever its size (the rainfall is given \emph{per unit area}, so the size cancels). The answer is $\boxed{\textbf{(D)}\ 1.5\text{ cm}}$, not (E).
